\section{生成模型的统一框架：从 Latent 到 Data}

\subsection{核心思想：学习分布间的映射}

\begin{importantbox}{所有生成模型的共同思路}
所有深度生成模型（GAN、VAE、Diffusion Model、Flow Model）共享同一个核心思想：\textbf{学习一个从简单分布（latent distribution）到复杂数据分布（data distribution）的映射}。

具体而言：
\begin{enumerate}
    \item 选择一个容易采样的简单分布作为 latent distribution $p(z)$（通常是标准高斯分布 $\mathcal{N}(\mathbf{0}, \mathbf{I})$）
    \item 用神经网络学习一个映射 $G_\theta: z \mapsto x$
    \item 生成时：采样 $z \sim p(z)$，计算 $x = G_\theta(z)$
\end{enumerate}
不同模型的区别在于\textbf{如何训练这个映射}。
\end{importantbox}

\subsection{符号约定}

本课程使用以下贯穿始终的符号：

\begin{table}[H]
\centering
\begin{tabular}{lll}
\toprule
符号 & 含义 & 说明 \\
\midrule
$z$ & latent 变量 & 来自简单分布（通常为 $\mathcal{N}(\mathbf{0}, \mathbf{I})$） \\
$x$ & 数据点 & 来自训练数据集 \\
$p(z)$ & latent / prior 分布 & 已知 \\
$p(x)$ / $p_{\text{data}}(x)$ & 数据分布 & \textbf{未知}，这是我们要逼近的目标 \\
$p_\theta(x|z)$ & likelihood 分布 & 给定 $z$ 生成 $x$ 的条件分布 \\
$p(z|x)$ & posterior 分布 & 给定 $x$ 推断 $z$ 的条件分布 \\
\bottomrule
\end{tabular}
\caption{本课程使用的符号约定}
\end{table}

\subsection{Autoencoder 的启发}

Autoencoder 是理解生成模型的一个起点：

\begin{itemize}
    \item \textbf{Encoder}：$f_\phi: x \mapsto z$，将数据点压缩到低维 latent 空间
    \item \textbf{Decoder}：$g_\theta: z \mapsto \hat{x}$，从 latent 表示重建数据
    \item 训练目标：最小化重建误差 $\|x - g_\theta(f_\phi(x))\|^2$
\end{itemize}

在生成模型的语境下，\textbf{decoder 就是我们需要的映射}——它将 latent 空间中的点映射到数据空间。但关键问题是：\textbf{如何保证 decoder 能够将 latent distribution 中的任意采样点都映射为合理的数据点？}

\begin{warningbox}{Autoencoder 不等于生成模型}
普通 Autoencoder 的 latent 空间并不服从某个已知的简单分布。如果我们直接从 latent 空间随机采样一个点 $z$ 并送入 decoder，得到的输出很可能不像真实数据。VAE 正是通过对 latent 空间施加先验约束来解决这个问题。
\end{warningbox}

\subsection{本章小结}

生成模型的统一范式是：从一个已知的简单分布 $p(z)$ 出发，通过神经网络学习到数据分布 $p_{\text{data}}(x)$ 的映射。GAN、VAE、Diffusion Model 和 Flow Model 的区别仅在于训练策略的不同。Autoencoder 的 decoder 提供了这种映射的雏形，但需要额外的机制来保证 latent 空间的结构化。

\newpage
